\subsection{معماری و پیاده‌سازی}
پیاده‌سازی ماژول در فایل \lr{Codes/framework/compression\_engine/compot\_compressor.py} مستقر بوده و شامل ارکان زیر است:
\begin{itemize}
    \item \textbf{کلاس \lr{COMPOTLinear}:} لایه فاکتورسازی‌شده جایگزین \lr{nn.Linear} و \lr{Conv1D} که عملگر گذر رو به جلو را در دو مرحله زیر-مربعی اجرا می‌کند:
    \begin{equation}
    y = ((x \cdot A^\top) \cdot B^\top) + b
    \end{equation}
    این کلاس ویژگی \lr{is\_custom\_attention = True} را تنظیم کرده و از بازسازی کامل وزن برحسب نیاز با متد \lr{get\_weight()} پشتیبانی می‌نماید.
    \item \textbf{متد کلاسی \lr{from\_linear\_and\_activations}:} کالبدشکافی وزن لایه اصلی، بررسی جهت ترانهاده برای لایه‌های خاص \lr{Conv1D} با ابعاد معکوس، استخراج فعال‌سازی‌ها، محاسبه خطای بازسازی نسبی (\lr{reconstruction\_error}) و برگرداندن لایه فشرده با رتبه مهارشده.
    \item \textbf{کلاس \lr{COMPOTCompressor}:} مدیریت کشف لایه‌های هدف ترنسفورمر، فیلتر کردن لایه‌های نامناسب طبق سازوکار \lr{should\_compress\_linear}، جمع‌آوری تانسورهای فعال‌سازی کالیبراسیون و جایگزینی ماژولار در گراف مدل.
\end{itemize}

\subsection{نتایج تجربی}
ارزیابی آزمایشگاهی ماژول \lr{compot\_compressor} بر روی مجموعه‌داده \lr{GLUE SST-2} و بنچمارک استاندارد رساله در جدول \ref{tab:benchmark_compot_compressor} خلاصه شده است:

\begin{table}[H]
\centering
\small
\caption{نتایج بنچ‌مارک تجربی ماژول \lr{compot\_compressor} روی مدل‌های پایه}
\label{tab:benchmark_compot_compressor}
\begin{tabular}{lccccc}
\toprule
\textbf{مدل پایه \lr{(Base Model)}} & \textbf{دقت \lr{(Accuracy)}} & \textbf{شاخص \lr{F1}} & \textbf{حافظه \lr{(Peak VRAM)}} & \textbf{توان \lr{(Throughput)}} & \textbf{وضعیت} \\
\midrule
\lr{DistilBERT} & \lr{75.92\%} & \lr{0.8155} & \lr{368.29 MB} & \lr{64,689.20 tok/s} & \lr{PASS} \\
\lr{RoBERTa} & \lr{50.92\%} & \lr{0.6748} & \lr{514.29 MB} & \lr{12,361.10 tok/s} & \lr{SKIP\_ROBERTA\_ARCHITECTURAL} \\
\lr{GPT-2} & \lr{52.18\%} & \lr{0.6854} & \lr{576.21 MB} & \lr{1,643.37 tok/s} & \lr{SKIP\_GPT\_ARCHITECTURAL} \\
\lr{Qwen} & -- & -- & -- & -- & \lr{NOT\_RUN} \\
\bottomrule
\end{tabular}
\end{table}

ماژول بر روی مدل \lr{DistilBERT} با موفقیت به دقت ۷۵٫۹۲٪ در غربالگری سریع و ۹۰٫۸۳٪ در جاروب رتبه‌ای (با رتبه ۶۴ و شاخص F1 معادل ۰٫۹۱۰۵) دست یافته است. با این حال، در مدل‌های \lr{RoBERTa} و \lr{GPT-2} در تنظیمات کالیبراسیون بدون آداپتور ریکاوری عمیق، دقت در فلات کلاسی ۵۰٫۹۲٪ متوقف گردیده و برچسب رد معماری دریافت کرده است.

\subsection{محدودیت‌ها و تحلیل شکست}
بررسی لاگ‌های خطا و رفتار کلاسی ماژول نشان‌دهنده چالش‌های فنی زیر است:
\begin{itemize}
    \item \textbf{افتادن در فلات شانس کلاسی در مدل‌های بدون بازآموزی:} در ارزیابی مدل‌های \lr{RoBERTa} و \lr{GPT-2}، دقت اولیه در حد اکثریت کلاس داده ($50.92\% \le 56.00\%$) ثبت گردید که موجب صدور وضعیت‌های \lr{SKIP\_ROBERTA\_ARCHITECTURAL} و \lr{SKIP\_GPT\_ARCHITECTURAL} گردید. عدم بازآموزی پس از فاکتورسازی تجزیه \lr{QR} بدون لایه \lr{LoRA} فعال، به ناپایداری فضای پنهان لایه‌های عمیق می‌انجامد.
    \item \textbf{خطای ناپایداری شبه‌وارون \lr{pinv} در ماتریس‌های کم‌انرژی:} در ماتریس‌های فعال‌سازی با رتبه ناکامل، ماتریس $R_X$ دارای مقادیر قطری نزدیک به صفر می‌گردد که محاسبه $(R_X^\top)^\dagger$ را با خطای عددی شدید یا تولید مقادیر \lr{NaN/Inf} مواجه می‌سازد؛ اگرچه پیاده‌سازی سورس با بلوک \lr{try-except} این خطا را پوشش داده و به مقداردهی تصادفی بازمی‌گردد، اما این عقب‌نشینی به بهای از دست رفتن انطباق داده‌محور تمام می‌شود.
\end{itemize}
